\begin{center}
    \Large{\textbf{Аннотация}}
\end{center}

Декодирование визуальных стимулов по нейробиологическим сигналам представляет собой ключевую задачу на стыке вычислительной нейронауки и современного машинного обучения. Несмотря на успехи, достигнутые в одномодальных подходах с применением контрастивных представлений и диффузионных генеративных моделей, большинство существующих методов опирается лишь на одну модальность~--- либо функциональную магнитно-резонансную томографию (фМРТ) с высоким пространственным разрешением, либо электроэнцефалографию (ЭЭГ) с высоким временным разрешением. Совместное использование этих модальностей остаётся слабо изученным направлением.

В работе предлагается мультимодальная архитектура, одновременно обрабатывающая записи фМРТ-ЭЭГ для реконструкции визуальных стимулов. Эмбеддинги нейробиологических сигналов обучаются контрастивно для сближения с CLIP-эмбеддингами соответствующих изображений. Двухстадийная схема генерации включает уточнение эмбеддинга мозга в CLIP-пространстве априорной диффузионной моделью и последующую генерацию изображения предобученной диффузионной моделью. Гемодинамическая задержка BOLD-сигнала анализируется через решение вспомогательной обратной задачи прогнозирования фМРТ по визуальному стимулу; полученная оценка задержки используется на этапе предобработки при формировании обучающих триплетов.

Эксперименты на открытом мультимодальном наборе данных подтверждают эффективность предложенной архитектуры в задаче нейродекодирования. По метрике CLIP-Score мультимодальная модель превосходит одномодальные аналоги, что подчёркивает необходимость совместного учёта фМРТ и ЭЭГ для точной реконструкции визуальных стимулов.
